\documentclass[11pt,a4paper]{article}
\usepackage[T1]{fontenc}
\usepackage[utf8]{inputenc}
\usepackage{amsmath,amsthm,mathtools}
\usepackage{newtxtext,newtxmath}
\usepackage[a4paper,margin=27mm,headheight=15pt]{geometry}
\usepackage{microtype}

\usepackage{booktabs,array,longtable}
\usepackage[dvipsnames]{xcolor}
\usepackage[most]{tcolorbox}
\usepackage{enumitem,needspace}
\usepackage{fancyhdr}
\usepackage{titlesec}
\usepackage{listings}
\usepackage{hyperref}
\usepackage{bookmark}
\definecolor{Ink}{HTML}{19384A}
\definecolor{Teal}{HTML}{006D77}
\definecolor{Pale}{HTML}{F0F6F7}
\definecolor{Muted}{HTML}{52626D}
\hypersetup{colorlinks=true,linkcolor=Teal,citecolor=Teal,urlcolor=Teal,
  pdftitle={Friendly Order Type Through Incomparability Components},
  pdfauthor={ChatGPT},pdfsubject={Finite posets, ordinal invariants, and multiset orders}}
\titleformat{\section}{\Large\bfseries\color{Ink}}{\thesection}{0.7em}{}
\titleformat{\subsection}{\large\bfseries\color{Ink}}{\thesubsection}{0.7em}{}
\pagestyle{fancy}
\fancyhf{}
\fancyhead[L]{\small\color{Muted}Friendly order type and incomparability}
\fancyhead[R]{\small\color{Muted}Research note}
\fancyfoot[C]{\small\thepage}
\renewcommand{\headrulewidth}{0.3pt}
\setlength{\parskip}{0.35em}
\setlength{\parindent}{1.2em}
\setlength{\emergencystretch}{2em}
\setcounter{tocdepth}{2}
\newtheorem{theorem}{Theorem}[section]
\newtheorem{lemma}[theorem]{Lemma}
\newtheorem{proposition}[theorem]{Proposition}
\newtheorem{corollary}[theorem]{Corollary}
\theoremstyle{definition}
\newtheorem{definition}[theorem]{Definition}
\newtheorem{example}[theorem]{Example}
\newtheorem{problem}[theorem]{Problem}
\theoremstyle{remark}
\newtheorem{remark}[theorem]{Remark}
\newcommand{\fot}{\operatorname{o}_{\perp}}
\newcommand{\mot}{\operatorname{o}}
\newcommand{\wid}{\operatorname{w}}
\newcommand{\hei}{\operatorname{h}}
\newcommand{\incg}{G_{\perp}}
\newcommand{\cc}{c_{\perp}}
\newcommand{\Bad}{\operatorname{Bad}}
\newcommand{\Fr}{\operatorname{Bad}_{\perp}}
\newcommand{\Mr}{\mathsf{M}^{r}}
\newcommand{\rk}{\operatorname{rk}}
\newcommand{\pred}{\operatorname{pred}}
\newcommand{\N}{\mathbb{N}}
\newcommand{\F}{\mathbb{F}}
\newcommand{\res}[2]{#1_{\not\ge #2}}
\newcommand{\CA}{\mathsf{C}}
\newcommand{\Ant}{\mathsf{A}}
\newcommand{\bool}{\mathsf{B}}
\newtcolorbox{resultbox}{colback=Pale,colframe=Teal,boxrule=0.7pt,
  arc=1mm,left=4mm,right=4mm,top=3mm,bottom=3mm}
\lstset{basicstyle=\small\ttfamily,keywordstyle=\color{Teal}\bfseries,
  commentstyle=\color{Muted},breaklines=true,columns=fullflexible,
  keepspaces=true,showstringspaces=false,frame=single,rulecolor=\color{Teal},
  xleftmargin=0pt,xrightmargin=0pt,aboveskip=1em,belowskip=1em}

\begin{document}
\thispagestyle{empty}
\noindent{\small\sffamily\bfseries\color{Teal}ORDER THEORY\quad /\quad ORDINAL INVARIANTS\quad /\quad 19 SEPTEMBER 2026}
\vspace{1.1cm}

\noindent{\Huge\bfseries\color{Ink}Friendly Order Type\\[0.12em]
Through Incomparability\\[0.12em]Components}
\vspace{0.5cm}

\noindent{\Large A finite structural case of an open problem,\\
with Cartesian products and transfinite extensions}
\vspace{0.7cm}

\begin{resultbox}
For every finite poset $P$,
\[
\boxed{\fot(P)=|P|-\cc(P)}
\qquad\text{and hence}\qquad
\boxed{\wid(\Mr(P))=\omega^{\,|P|-\cc(P)}}.
\]
Here $\cc(P)$ counts the connected components of the incomparability graph,
including isolated vertices. The first quantity is also the rank of that
graph's graphic matroid.
\end{resultbox}

\subsection*{Abstract}
Vialard introduced friendly order type to compute the ordinal width of the
Dershowitz--Manna multiset ordering and asked for further structural and
compositional descriptions. We study the finite structural case. We prove
that the friendly order type of a finite poset is its number of elements
minus its number of incomparability components. One proof bounds every
open-ended bad sequence by a graph-forest potential; the matching lower
bound constructs connected linear extensions of the components. This
identifies the invariant with a familiar graph rank and gives polynomial-time
value and witness algorithms.

We derive exact formulas for finite lexicographic substitutions and Cartesian
products. We also prove additivity over arbitrary well-ordered ordinal sums,
which yields a transfinite evaluation theorem for all well-partial-orders
with finite incomparability components. Explicit countable examples show
that neither graph-isomorphism invariance nor a simple connected-component
correction to maximal order type extends to arbitrary infinite bases.
Independent residual-rank computations verify the finite formula on all
101,660 naturally labelled posets of sizes at most seven; additional product
and substitution checks are supplied as executable Python.

\vspace{0.3cm}
\noindent\textbf{Research status.}
The proofs below are complete mathematical arguments, but have not been
externally refereed or proof-assistant checked. The result solves a specified
finite subproblem and an infinite extension, not the entire published
research program. A targeted literature search did not locate the component
formula; that search does not establish priority. The general multiset-width
theorem is Vialard's, not a new claim of this manuscript.

\newpage
\begingroup
\small
\setlength{\parskip}{0pt}
\tableofcontents
\endgroup
\newpage

\section{The documented problem and the scope of the answer}

In the conclusion of \emph{Ordinal Measures of the Set of Finite Multisets},
Vialard asks about the friendly order type: ``How does it relate to other
concepts? Can it be computed compositionally for more operations?''
\cite[p.~87:11]{vialard2023}.
The 2024 doctoral thesis again asks for compositional computation and a
deeper interpretation of this invariant \cite[Conclusion, pp.~109--110]{vialard2024}.
These are open-ended research questions rather than one numbered yes/no
conjecture.

The concrete problem selected here is the following finite structural
specialization of that program.

\begin{problem}[Finite structural subproblem]\label{prob:finite}
Find an exact structural expression for $\fot(P)$ for every finite poset $P$,
relate it to a familiar invariant, and use that expression to compute it for
additional poset constructions. Determine how far the expression extends
to infinite well-partial-orders.
\end{problem}

The literature consulted does not separately designate
Problem~\ref{prob:finite} as a named open conjecture. We therefore do not
claim to have settled a previously named conjecture. The contribution is a
proposed structural solution to a concrete part of the documented program.
The finite case is not merely a question about finite ordinal widths:
$\Mr(P)$ is usually infinite even when $P$ is finite, and its width is the
transfinite ordinal $\omega^{\fot(P)}$.

The principal statements are as follows.

\begin{center}
\renewcommand{\arraystretch}{1.2}
\begin{tabular}{@{}p{0.27\textwidth}p{0.65\textwidth}@{}}
\toprule
Domain & Exact conclusion \\
\midrule
Finite $P$ & $\fot(P)=|P|-\cc(P)$; equivalently, graphic-matroid rank. \\
Finite substitutions & Each nontrivial index component merges all its fibers into one incomparability component. \\
Finite $A\times B$, $|A|,|B|\ge2$ & $\fot(A\times B)=|A||B|-1-b-t$, where $b,t$ indicate the existence of a least and greatest element in the product. \\
Well-ordered sums & $\fot(\sum_{i<\delta}P_i)=\sum_{i<\delta}\fot(P_i)$, using ordinary ordinal addition. \\
Finite-component wpos & The preceding sum becomes $\sum_{i<\delta}(|C_i|-1)$ over the ordered incomparability components. \\
\bottomrule
\end{tabular}
\end{center}

The finite component theorem, its algorithmic form, the substitution formula,
and the Cartesian formula are proved independently of the multiset-width
theorem. The transfinite argument uses only the residual definition of
friendly order type and well-founded tree ranks. The conversion from friendly
order type to multiset width invokes Vialard's theorem explicitly.

\section{Definitions and conventions}

\subsection{Posets, residuals, and tree ranks}

All orders in this article are partial orders, not merely quasi-orders.
A \emph{well-partial-order} (wpo) is a poset with no infinite bad sequence.
A finite sequence $x_0,\ldots,x_{k-1}$ in $P$ is \emph{bad} when
\[
 i<j<k\quad\Longrightarrow\quad x_i\not\le_P x_j.
\]
Every finite poset is a wpo. For $x\in P$, set
\[
 \uparrow_P x=\{y\in P:x\le_P y\},\qquad
 \res{P}{x}=P\setminus\uparrow_P x.
\]
All subsets carry the induced order. After a bad sequence $s$, the residual
is $P\setminus\bigcup_{x\in s}\uparrow_P x$. This is precisely the set of
points that can still be appended to $s$ while preserving badness.

For a well-founded rooted tree, a leaf has rank zero, and a node has rank
\[
 \sup\{\text{rank of a child}+1:\text{child of the node}\}.
\]
The rank of the tree is the rank of its root. An empty supremum is zero.
In particular, the tree consisting only of its empty sequence has rank zero.
The tree $\Bad(P)$ of finite bad sequences is well-founded for a wpo.
We write $\mot(P)$ for its rank. This is the usual maximal-order-type
invariant, but only its rank definition is needed in our transfinite proof.
It satisfies
\begin{equation}\label{eq:mot}
 \mot(P)=\sup_{x\in P}\bigl(\mot(\res{P}{x})+1\bigr).
\end{equation}
For a finite poset, $\mot(P)=|P|$: reverse any linear extension to obtain a
bad sequence containing all elements.

\subsection{Friendly order type}

We use Vialard's definition of open-ended bad sequences
\cite[Definition~5.3.2]{vialard2024}.
Two points are \emph{incomparable}, written $x\perp y$, when neither
$x\le y$ nor $y\le x$ holds. Write
$P_{\perp x}=\{y\in P:y\perp x\}$.

\begin{definition}[Open-ended sequence and friendly order type]
The empty sequence is open-ended. If $s$ is open-ended and its current
residual is $R$, then $sx$ is open-ended when $x\in R$ has an incomparable
point $y\in R$. Such a point $y$ is called a \emph{friend} of $x$ at that
step. Let $\Fr(P)$ be the tree of open-ended bad sequences. Its rank is the
\emph{friendly order type} $\fot(P)$.
\end{definition}

The friend need not be a member of the eventual selected sequence, and need
not remain forever. It must exist in the residual \emph{at the step where
it is used}. Incomparability guarantees that it is not deleted by selecting
$x$, because $x\not\le y$.

The defining recurrence is
\begin{equation}\label{eq:fot}
 \fot(P)=\sup_{\substack{x\in P\\P_{\perp x}\ne\varnothing}}
 \bigl(\fot(\res{P}{x})+1\bigr).
\end{equation}
For finite $P$, this is a maximum, with value zero if no move is permitted,
and the rank equals the maximum length of an open-ended bad sequence.
In particular, $\fot(\varnothing)=0$ and $\fot(P)=0$ for every chain.

\subsection{The graph and the multiset order}

The \emph{incomparability graph} $\incg(P)$ is the simple undirected graph
with vertex set $P$ and edges $\{x,y\}$ exactly when $x\perp y$.
For finite $P$, let $\cc(P)$ be its number of connected components.
Isolated vertices count as components, and $\cc(\varnothing)=0$.
For any finite graph $G$, write
\[
 r(G)=|V(G)|-c(G).
\]
We will prove directly that this is both the maximum number of edges in a
forest contained in $G$ and the rank of an incidence matrix.

For completeness, the multiset ordering used here is not multiset embedding.
For finite multisets $m,n$ over $P$, cancel their common multiplicities and
write $m'=m\setminus(m\cap n)$ and $n'=n\setminus(m\cap n)$. Then
\begin{equation}\label{eq:multiset}
 m\le_r n\quad\Longleftrightarrow\quad
 m=n\ \text{or}\ (\forall x\in m')\,(\exists y\in n')\ x<_P y.
\end{equation}
No injective matching of the remaining occurrences is required. Write
$\Mr(P)$ for the resulting ordered set of finite multisets, including the
empty multiset \cite[Definition~1.6]{vialard2023}.

The ordinal width $\wid(Q)$ of a wpo $Q$ is the rank of its tree of finite
sequences of pairwise incomparable points. For finite $Q$ it equals the
ordinary maximum antichain size. For infinite $Q$ it is generally not the
supremum of the finite antichain sizes: the rank also records the nested
structure of the possible continuations.

\begin{theorem}[Vialard]\label{thm:vialard}
For every wpo $P$,
\[
 \wid(\Mr(P))=\omega^{\fot(P)}.
\]
\end{theorem}

This is \cite[Theorem~3.4]{vialard2023}. We use it as a published theorem;
we do not claim to reprove its general multiset argument. If $P$ is empty,
$\Mr(P)$ has one element and the formula correctly gives $1=\omega^0$.

\section{The finite component theorem}

\begin{theorem}[Finite component formula]\label{thm:finite}
For every finite poset $P$,
\[
 \boxed{\fot(P)=|P|-\cc(P).}
\]
Moreover, an open-ended bad sequence of this length can be constructed
without evaluating the residual-rank recurrence.
\end{theorem}

The proof separates a graph-theoretic upper bound from an order-theoretic
construction. The latter separation matters: connectivity alone does not
make an arbitrary vertex ordering into a bad sequence.

\subsection{Upper bound: every legal move consumes forest rank}

\begin{lemma}[Graph-rank decrease]\label{lem:potential}
Let $R$ be a finite poset, let $x\in R$ have a friend $y\in R$, and put
$S=R\setminus\uparrow_R x$. Then
\[
 r(\incg(S))\le r(\incg(R))-1.
\]
\end{lemma}

\begin{proof}
A forest on a graph with $v$ vertices and $c$ connected components has at
most $v-c$ edges: in each component, an acyclic graph on its vertices has at
most one fewer edge than vertices. Choosing a spanning tree in each
component attains $v-c$. Thus $r(G)$ is the maximum forest size.

Choose a spanning forest $F$ of $\incg(S)$. Since $x\perp y$, the point $y$
belongs to $S$, whereas $x$ does not. Adding the new vertex $x$ and the edge
$\{x,y\}$ to $F$ cannot create a cycle. It produces a forest in $\incg(R)$
with $r(\incg(S))+1$ edges. The asserted inequality follows.
\end{proof}

Every step of an open-ended sequence decreases this nonnegative integer
potential by at least one. Therefore
\begin{equation}\label{eq:upper}
 \fot(P)\le r(\incg(P))=|P|-\cc(P).
\end{equation}
This proof allows an upset deletion to remove many points and to split old
components into new ones. It does not assume that a move simply removes one
vertex, or that connectivity is preserved after every possible move.

\subsection{Incomparability components are ordinal-sum blocks}

\begin{lemma}[Ordered components]\label{lem:components}
Distinct connected components of $\incg(P)$ are uniformly comparable:
for two such components $C,D$, either every point of $C$ is below every
point of $D$, or every point of $D$ is below every point of $C$.
These comparisons linearly order the components. This holds for arbitrary
posets, finite or infinite.
\end{lemma}

\begin{proof}
Points in different components must be comparable. Fix a point $z$ outside a
component $C$. If $u\perp v$ are adjacent vertices in $C$, the directions of
the comparisons with $z$ cannot differ: $u<z<v$ would imply $u<v$, and
$v<z<u$ would imply $v<u$. Following a finite graph path shows that either
all of $C$ is below $z$ or all of $C$ is above $z$.

Now vary $z$ along a path in a second component $D$; the same argument shows
that the direction is constant throughout $D$. This proves uniformity.
Totality follows from comparability between components, and transitivity
follows from transitivity in $P$.
\end{proof}

For finite $P$ we may consequently write
\[
 P=C_1+C_2+\cdots+C_k,
\]
where $+$ is ordinal sum: every point of an earlier block is below every
point of a later block. Each induced graph $\incg(C_i)$ is connected.

\subsection{Connected linear extensions}

\begin{lemma}[A boundary point with a friend]\label{lem:boundary}
Let $C$ be a finite poset whose incomparability graph is connected. If $D$
is a nonempty proper order ideal of $C$, then some minimal point $v$ of
$C\setminus D$ is incomparable to a point of $D$.
\end{lemma}

\begin{proof}
Suppose the assertion fails. Every minimal point $v$ of $C\setminus D$ is
then comparable to every $d\in D$. We cannot have $v\le d$, because $D$
is an order ideal and $v\notin D$. Thus $d<v$ for every such $d,v$.

Every point $z$ of the finite poset $C\setminus D$ lies above a minimal
point $v$ of that poset. Hence $d<v\le z$ for every $d\in D$. There are
therefore no incomparability edges between the two nonempty sets $D$ and
$C\setminus D$, contradicting connectivity.
\end{proof}

\begin{lemma}[Connected-prefix extension]\label{lem:extension}
If $C$ is a finite poset with connected incomparability graph and $|C|=m$,
there is a linear extension $v_1,\ldots,v_m$ such that every $v_j$ with
$j\ge2$ has an incomparable point among $v_1,\ldots,v_{j-1}$.
In particular, every nonempty prefix induces a connected incomparability
graph.
\end{lemma}

\begin{proof}
Start with any minimal point $v_1$. Suppose an ideal
$D=\{v_1,\ldots,v_j\}$ has been constructed. If it is proper,
Lemma~\ref{lem:boundary} supplies a minimal remaining point $v_{j+1}$ with
a friend in $D$. Appending a minimal remaining point preserves the
linear-extension property and keeps the selected set an ideal. Its friend
connects the new point to the existing prefix. Continue until all points
have been selected. The one-point case is included.
\end{proof}

\begin{corollary}\label{cor:connected}
If $\incg(C)$ is connected and $C$ is finite, then
$\fot(C)=|C|-1$.
\end{corollary}

\begin{proof}
Use Lemma~\ref{lem:extension} and read the extension backwards, omitting its
first point:
\[
 v_m,v_{m-1},\ldots,v_2.
\]
After the later points have been removed, $v_j$ is maximal in the remaining
prefix $\{v_1,\ldots,v_j\}$. Its upset in that prefix is therefore the
singleton $\{v_j\}$. It also has an incomparable friend in the earlier
prefix. Every move is legal, giving length $m-1$. The reverse inequality is
\eqref{eq:upper}.
\end{proof}

\subsection{Completion of the proof and an explicit certificate}

\begin{proof}[Proof of Theorem~\ref{thm:finite}]
The empty case is immediate. Decompose a nonempty $P$ into its ordered
incomparability components $C_1<\cdots<C_k$. For each $C_i$, choose the
extension of Lemma~\ref{lem:extension}, and retain its reverse with the first
point omitted. Process these lists in the order $C_k,C_{k-1},\ldots,C_1$.

While working in $C_i$, all points of earlier components are still present:
selecting a point in a later component never removes a smaller point. The
unprocessed prefix inside $C_i$ is also still present. A chosen point has
its designated friend in that prefix, and is maximal within that prefix.
Its upset may delete leftover points of higher components, but this cannot
affect its friend or any future move in a lower component. Thus the entire
concatenated sequence is open-ended.

Its length is
\[
 \sum_{i=1}^k(|C_i|-1)=|P|-k=|P|-\cc(P).
\]
Together with \eqref{eq:upper}, this proves the theorem.
\end{proof}

Choose, for each $v_j$ in a component extension with $j\ge2$, one earlier
friend $p(v_j)$. The edges $\{v_j,p(v_j)\}$ form a spanning tree of the
component: each edge points toward an earlier index, and repeated parent
steps reach $v_1$. Over all components they form a spanning forest with
$|P|-\cc(P)$ edges. The same data therefore certify both a maximal forest
and an optimal open-ended sequence.

Exactly one point of each component is omitted from the \emph{selected
sequence}. This does not mean that one point of each original component
survives in the final residual. A later move in a lower component can delete
an omitted point of a higher component. This distinction is important for
checking the construction.

\section{Graph rank, consequences, and small examples}

\subsection{A familiar invariant, not an identification of structures}

\begin{corollary}[Graphic-matroid and incidence rank]\label{cor:matroid}
For a finite poset $P$ and any field $K$,
\[
 \fot(P)=r(\incg(P))=\rk_K B,
\]
where $B$ is any oriented vertex--edge incidence matrix of $\incg(P)$.
Equivalently, $\fot(P)$ is the rank of the graphic matroid of $\incg(P)$.
\end{corollary}

\begin{proof}
The forest characterization of $r(G)$ was established in
Lemma~\ref{lem:potential}. For the matrix statement, each incidence column
has coordinate sum zero within its connected component. Hence the column
rank is at most $|V(G)|-c(G)$. The columns belonging to a spanning forest
are independent over every field: a leaf row singles out its incident edge,
whose coefficient must vanish in any dependence, and leaf removal finishes
the induction. The lower bound is therefore the same number.
\end{proof}

In characteristic two, the signs disappear, so an edge column simply has
a $1$ at each endpoint. This provides an independent, easily implemented
algebraic check. The corollary identifies \emph{ranks}; it does not identify
the entire tree of friendly sequences with a graphic matroid, and does not
assert that every forest orientation gives a legal sequence.

\Needspace{11\baselineskip}
\subsection{Monotonicity and one-point changes}

\begin{corollary}\label{cor:consequences}
For finite posets the following statements hold.
\begin{enumerate}[label=(\roman*),itemsep=0.2em]
\item $\fot(P)=\fot(P^{\mathrm{op}})$, and isomorphic incomparability graphs
have equal friendly order type.
\item $\fot(P)=0$ exactly when $P$ is a chain.
\item For nonempty $P$, $\fot(P)\le |P|-1$, with equality exactly when its
incomparability graph is connected.
\item Passing to an induced subposet, or adding order comparisons while
retaining a partial order, cannot increase friendly order type.
\end{enumerate}
\end{corollary}

\begin{proof}
Duality preserves incomparability, proving (i). Graph rank is zero precisely
for an edgeless graph, proving (ii). Statement (iii) is the component formula.
For (iv), deleting vertices or edges cannot increase the maximum size of a
forest. Adding comparisons deletes incomparability edges.
\end{proof}

\begin{proposition}[Insertion rule]\label{prop:insertion}
Let $P'$ be a finite poset obtained from an induced subposet $P$ by adding
one point $x$. Let $t$ be the number of components of $\incg(P)$ that contain
at least one friend of $x$. Then
\[
 \fot(P')-\fot(P)=t.
\]
\end{proposition}

\begin{proof}
If $t=0$, the new point is isolated and both the vertex count and component
count increase by one. If $t>0$, the new point merges exactly those $t$
components into one, so $c(\incg(P'))=c(\incg(P))-t+1$. Substitute into
Theorem~\ref{thm:finite}.
\end{proof}

\subsection{Chains, antichains, and equal ordinary invariants}

Write $\CA_n$ for an $n$-element chain and $\Ant_n$ for an $n$-element
antichain. For $n\ge1$,
\[
 \fot(\CA_n)=0,\qquad \fot(\Ant_n)=n-1.
\]
For nonempty finite $A,B$, the ordinal sum and disjoint sum satisfy
\[
 \fot(A+B)=\fot(A)+\fot(B),\qquad
 \fot(A\sqcup B)=|A|+|B|-1.
\]
The first follows because there are no cross-block incomparability edges;
the second because every cross-block pair is incomparable, making the
graph connected. These two special operation rules are consistent with
Vialard's existing binary-sum formulas; they are not claimed here as new
binary-sum results \cite[Proposition~4.1]{vialard2023}.

\begin{example}[A four-point separation]\label{ex:four}
Let
\[
 P=\Ant_2+\Ant_2,\qquad Q=\CA_2\sqcup\CA_2.
\]
Both have four elements, height two, ordinary width two, and maximal order
type four. The incomparability graph of $P$ consists of two disjoint edges,
whereas the graph of $Q$ is the connected graph $K_{2,2}$. Thus
\[
 \fot(P)=2,\qquad \fot(Q)=3,
\]
and Theorem~\ref{thm:vialard} gives
\[
 \wid(\Mr(P))=\omega^2,\qquad \wid(\Mr(Q))=\omega^3.
\]
This directly demonstrates, already on finite bases, that the usual three
ordinal invariants of a base do not determine its multiset-order width.
\end{example}

For every fixed $n\ge1$, every integer $r$ with $0\le r\le n-1$ occurs as
$\fot(P)$ on an $n$-element poset. One realization is
\[
 P=\Ant_{r+1}+\CA_{n-r-1},
\]
where a zero-sized final chain is omitted. Consequently the possible
multiset-order widths of $n$-element bases are exactly
$1,\omega,\ldots,\omega^{n-1}$.

\section{Finite lexicographic substitutions}

\begin{definition}[Substitution]
Let $R$ be a finite poset and let $P_i$ be a nonempty finite poset for every
$i\in R$. The substitution $R[P_i:i\in R]$ has elements $(i,x)$ with
$x\in P_i$. Within a fiber it uses the order of $P_i$, and between distinct
fibers it declares every point of $P_i$ below every point of $P_j$ exactly
when $i<_R j$.
\end{definition}

Let $\mathcal C(R)$ denote the set of incomparability components of $R$.
A singleton component is an isolated vertex of $\incg(R)$.

\begin{theorem}[Substitution formula]\label{thm:substitution}
Under the preceding hypotheses,
\begin{align}
 \fot(R[P_i:i\in R])
 &=\sum_{\substack{D\in\mathcal C(R)\\|D|>1}}
       \left(\sum_{i\in D}|P_i|-1\right)
   +\sum_{\substack{\{i\}\in\mathcal C(R)}}\fot(P_i).
 \label{eq:substitution}
\end{align}
All sums in this formula are finite sums of integers.
\end{theorem}

\begin{proof}
Whenever $i\perp_R j$, all vertices of the two fibers are mutually joined
by cross-fiber incomparability edges. Suppose $D$ is a nontrivial connected
component of $\incg(R)$. Every vertex $i$ of $D$ has a neighbor in $D$.
Two vertices in the same fiber over $i$ can therefore be joined by a
length-two path through any vertex of a neighboring fiber. Vertices in
different fibers over $D$ can be joined by lifting a path in $D$. Hence the
union of the fibers over $D$ forms one connected component.

There are no edges between fibers lying over distinct index components.
For an isolated index $i$, the fiber retains exactly its own internal
components. Thus a nontrivial $D$ contributes
$\sum_{i\in D}|P_i|-1$ to graph rank, and an isolated index contributes
$|P_i|-\cc(P_i)=\fot(P_i)$. Apply Theorem~\ref{thm:finite}.
\end{proof}

\begin{corollary}[Constant fiber]\label{cor:constantfiber}
Suppose every fiber is the same nonempty finite poset $P$. If $R$ has $s$
isolated vertices and $q$ nontrivial incomparability components, then
\[
 \cc(R[P])=q+s\cc(P),\qquad
 \fot(R[P])=|R||P|-q-s\cc(P).
\]
\end{corollary}

This provides a genuinely structural compositional rule: an isolated index
requires the internal friendly order type of its fiber, while a nontrivial
index component forgets the fiber orders and retains only their sizes.
Empty fibers are excluded because they can destroy connecting paths; they
can be handled by first deleting their indices and applying the theorem
to the remaining index poset. An empty index yields the empty poset.

\section{Cartesian products: only the endpoints split off}

The Cartesian order is coordinatewise:
$(a,b)\le(a',b')$ exactly when $a\le a'$ and $b\le b'$.
The next structural lemma applies even without finiteness.

\begin{lemma}[No middle ordinal cut in a nontrivial product]\label{lem:cut}
Let $A,B$ be posets, each with at least two elements. There is no partition
$A\times B=L\mathbin{\dot\cup}U$ with $|L|,|U|\ge2$ and every point of $L$
strictly below every point of $U$.
\end{lemma}

\begin{proof}
Suppose there is such a partition. Let
$L_A,U_A\subseteq A$ and $L_B,U_B\subseteq B$ be its coordinate projections.
Every point of $L_A$ is at most every point of $U_A$, and similarly in $B$.
It follows by antisymmetry that
\[
 |L_A\cap U_A|\le1,\qquad |L_B\cap U_B|\le1.
\]
Also, $L_A\cup U_A=A$ and $L_B\cup U_B=B$.

If $a\in L_A\setminus U_A$, its entire row $\{a\}\times B$ lies in $L$.
Therefore there can be no $b\in U_B\setminus L_B$, whose entire column
would lie in $U$. We have proved
\[
 L_A\setminus U_A\ne\varnothing\ \Longrightarrow\ U_B\subseteq L_B.
\]
Reversing $L,U$ yields the analogous implication
$U_A\setminus L_A\ne\varnothing\Rightarrow L_B\subseteq U_B$.

Since $|L|\ge2$, one of its projections has at least two points; exchange
coordinates if necessary so that $|L_A|\ge2$. Then
$L_A\setminus U_A\ne\varnothing$, so $U_B\subseteq L_B$, making $U_B$ a
singleton. Since $|U|\ge2$, we now have $|U_A|\ge2$. Hence
$U_A\setminus L_A\ne\varnothing$, and $L_B\subseteq U_B$. Consequently
$B=L_B\cup U_B$ is a singleton, a contradiction.
\end{proof}

\begin{theorem}[Incomparability graph of a product]\label{thm:productgraph}
If $|A|,|B|\ge2$, the incomparability graph of $A\times B$ has exactly one
nontrivial connected component. Its only possible additional components
are the singleton least element and the singleton greatest element of the
product, when they exist.
\end{theorem}

\begin{proof}
The product is not a chain. If either factor has an incomparable pair,
fixing the other coordinate gives one in the product. Otherwise take
$a_0<a_1$ and $b_0<b_1$; then $(a_0,b_1)\perp(a_1,b_0)$.
Thus there is a nontrivial component $C$.

By Lemma~\ref{lem:components}, all other components are linearly ordered
below or above $C$. Let $L$ be the union of the components below $C$.
If $|L|\ge2$, then $L$ and its complement give a forbidden cut from
Lemma~\ref{lem:cut}, since the complement contains $C$. Thus $|L|\le1$.
Similarly the union of components above $C$ has at most one point.
If the lower point exists, it is the least point of the product; if the
upper point exists, it is the greatest. Conversely a least or greatest
point is comparable to every point and therefore isolated in the graph.
\end{proof}

For a nonempty poset $A$, let $b(A),t(A)\in\{0,1\}$ indicate the existence
of a least and greatest point, respectively.

\begin{corollary}[Finite Cartesian formula]\label{cor:product}
For finite $A,B$ with $|A|,|B|\ge2$,
\begin{equation}\label{eq:product}
 \boxed{\fot(A\times B)=|A||B|-1-b(A)b(B)-t(A)t(B).}
\end{equation}
Consequently
\[
 \wid(\Mr(A\times B))
   =\omega^{\,|A||B|-1-b(A)b(B)-t(A)t(B)}.
\]
\end{corollary}

\begin{proof}
The product has a least point exactly when both factors have least points,
and likewise for greatest points. The graph therefore has
$1+b(A)b(B)+t(A)t(B)$ components. Apply Theorems~\ref{thm:finite}
and~\ref{thm:vialard}.
\end{proof}

If either factor is empty, the product is empty and has friendly order type
zero. If one factor has one point, the product is isomorphic to the other
factor. These exceptions must be handled before using \eqref{eq:product}.

\begin{corollary}[Several factors and finite grids]
For $d\ge2$ finite factors $P_1,\ldots,P_d$, each of size at least two,
\[
 \fot\!\left(\prod_{i=1}^dP_i\right)
 =\prod_{i=1}^d|P_i|-1-\prod_{i=1}^db(P_i)-\prod_{i=1}^dt(P_i).
\]
In particular, for $m_i\ge2$,
\[
 \fot(\CA_{m_1}\times\cdots\times\CA_{m_d})=\prod_{i=1}^dm_i-3.
\]
For the Boolean lattice $\bool_d=\CA_2^d$ with $d\ge2$,
\[
 \fot(\bool_d)=2^d-3,\qquad
 \wid(\Mr(\bool_d))=\omega^{2^d-3}.
\]
\end{corollary}

The finite product formula is notably simpler than a general computation
of ordinary width: away from the empty and singleton exceptions, only
cardinalities and endpoint indicators are needed. For example,
$\fot(\CA_2\times\CA_3)=3$ and $\fot(\bool_3)=5$.

\section{Well-ordered sums and a transfinite evaluation theorem}

\subsection{Additivity beyond binary sums}

For an ordinal $\delta$ and wpos $P_i$ ($i<\delta$), their ordinal sum
$P=\sum_{i<\delta}P_i$ puts every point of an earlier block below every point
of a later block. Empty blocks are permitted. This sum is a wpo: in an
infinite bad sequence the block indices would have to be nonincreasing,
hence eventually constant, giving an infinite bad sequence in one block.

The sums of ordinals below are \emph{ordinary ordinal sums}, not natural
sums and not commutative sums. Set
\[
 S_0=0,\quad S_{i+1}=S_i+\fot(P_i),\quad
 S_\lambda=\sup_{i<\lambda}S_i\quad(\lambda\text{ a limit}).
\]

\begin{theorem}[Well-ordered-sum additivity]\label{thm:transfinite}
For every ordinal $\delta$ and family of wpos $(P_i)_{i<\delta}$,
\begin{equation}\label{eq:transfinite}
 \boxed{\fot\!\left(\sum_{i<\delta}P_i\right)
          =\sum_{i<\delta}\fot(P_i).}
\end{equation}
\end{theorem}

\begin{proof}
We prove the assertion for all such representations simultaneously, by
transfinite induction on $\mot(P)$, where $P=\sum_{i<\delta}P_i$.
If $P$ is empty, both sides vanish. A point $x\in P_i$ has an incomparable
friend in $P$ exactly when it has one in $P_i$. Moreover,
\begin{equation}\label{eq:sumresidual}
 \res{P}{x}
 =\left(\sum_{j<i}P_j\right)+\res{(P_i)}{x}.
\end{equation}
Indeed, all later blocks are deleted, all earlier blocks survive, and only
the upset of $x$ is deleted in its own block.

Equation~\eqref{eq:mot} gives
$\mot(\res{P}{x})<\mot(P)$, so the induction hypothesis applies to the
representation on the right of \eqref{eq:sumresidual}. It yields
\[
 \fot(\res{P}{x})=S_i+\fot(\res{(P_i)}{x}).
\]
Substitution into \eqref{eq:fot} gives
\begin{align*}
 \fot(P)
 &=\sup_{\substack{i<\delta\\x\in P_i,\ (P_i)_{\perp x}\ne\varnothing}}
     \left(S_i+\fot(\res{(P_i)}{x})+1\right)\\
 &=\sup_{\substack{i<\delta\\\fot(P_i)>0}}
     \left(S_i+\fot(P_i)\right).
\end{align*}
The second equality uses associativity and continuity of ordinal addition
in its right argument, together with \eqref{eq:fot} within $P_i$.
Blocks with friendly order type zero permit no move and contribute no
increase to the partial sums. The displayed supremum is therefore exactly
$S_\delta=\sum_{i<\delta}\fot(P_i)$, including the case where all blocks
have value zero. This completes the induction.
\end{proof}

The binary case is already a published fact
\cite[Theorem~5.4.8(1)]{vialard2024}. The proof above records the arbitrary
well-ordered version and, in particular, makes the limit step explicit.
No assumption of countability, a computable presentation, or a bound on
component sizes is imposed.

\subsection{All wpos with finite incomparability components}

\begin{lemma}
The incomparability components of a wpo, ordered as in
Lemma~\ref{lem:components}, form a well-order.
\end{lemma}

\begin{proof}
Take any nonempty collection of components. The union of their points is
a nonempty subposet of a wpo, hence has a minimal point $x$. The component
containing $x$ is least in the collection: a lower component would supply
a point strictly below $x$. Thus every nonempty collection has a least
component.
\end{proof}

\begin{theorem}[Finite-component transfinite formula]\label{thm:finitecomponents}
Let $P$ be a wpo whose incomparability components are all finite. Enumerate
them in order as $(C_i)_{i<\delta}$. Then
\begin{equation}\label{eq:finitecomponents}
 \boxed{\fot(P)=\sum_{i<\delta}(|C_i|-1),\qquad
 \wid(\Mr(P))=\omega^{\,\sum_{i<\delta}(|C_i|-1)}.}
\end{equation}
More generally, if $P=\sum_{i<\delta}Q_i$ with arbitrary finite blocks, then
\[
 \fot(P)=\sum_{i<\delta}\bigl(|Q_i|-\cc(Q_i)\bigr).
\]
\end{theorem}

\begin{proof}
Each $C_i$ has connected incomparability graph, so its friendly order type
is $|C_i|-1$. Apply Theorem~\ref{thm:transfinite}, followed by
Theorem~\ref{thm:vialard}. For arbitrary finite blocks apply
Theorem~\ref{thm:finite} to each block first.
\end{proof}

This is not the ambiguous expression ``$|P|$ minus the number of
components'' for an infinite cardinal $|P|$. It is an ordinal sum of
explicit finite ranks, in the actual order of the components. Its order
sensitivity is essential.

\subsection{Ordinal examples and realizability}

If $Q$ is finite and $r=\fot(Q)$, then for every ordinal $\alpha$,
\begin{equation}\label{eq:repeat}
 \fot\!\left(\sum_{i<\alpha}Q\right)=r\cdot\alpha.
\end{equation}
The finite number $r$ is the \emph{left} factor in this ordinary ordinal
product. For $r=2$ and $\alpha=\omega+3$, the answer is $\omega+6$, not
$\omega\cdot2+6$.

\begin{example}[A transfinite sum of Boolean lattices]
Let $Q=\bool_3$, so $r=5$, and take
$\alpha=\omega^2+\omega\cdot3+5$. Then
\[
 \fot\!\left(\sum_{i<\alpha}\bool_3\right)
   =5\cdot\alpha=\omega^2+\omega\cdot3+25.
\]
Its multiset-order width is
\[
 \omega^{\omega^2+\omega\cdot3+25}.
\]
This base has infinitely many finite incomparability components, including
the isolated endpoints inside each Boolean-lattice block. They contribute
zero locally; the connected six-point middle contributes five.
\end{example}

\begin{corollary}[Every ordinal is a friendly order type]
For every ordinal $\alpha$, the wpo
$P_\alpha=\sum_{i<\alpha}\Ant_2$ satisfies
\[
 \fot(P_\alpha)=\alpha,\qquad
 \wid(\Mr(P_\alpha))=\omega^\alpha.
\]
\end{corollary}

This applies to uncountable ordinals as well as countable ones. It is an
existence statement about set-sized orders, not an effective notation system
for arbitrary ordinals. We regard it as an immediate corollary, not as a
separate priority claim.

\section{Why the infinite generalization needs more information}

The finite graph formula might suggest two stronger extrapolations: that
friendly order type is always determined by the abstract incomparability
graph, or that a connected wpo has maximal order type with a fixed one-point
correction. Both extrapolations fail. We give direct residual calculations,
so the examples do not depend on an ordinal-subtraction convention.

For ordinals $\alpha,\beta$, let
\[
 g(\alpha,\beta)=\fot(\alpha\sqcup\beta),
\]
where each ordinal is regarded as a chain. If either coordinate is zero,
$g(\alpha,\beta)=0$. If both are positive, every point has a friend in the
other chain and
\begin{equation}\label{eq:g}
 g(\alpha,\beta)=\sup\left(
  \{g(\xi,\beta)+1:\xi<\alpha\}
  \cup\{g(\alpha,\eta)+1:\eta<\beta\}\right).
\end{equation}
In particular, induction gives
\begin{align}
 g(m,n)&=m+n-1 &&(m,n\ge1\text{ finite}),\label{eq:gfinite}\\
 g(\alpha,1)&=\alpha &&(\alpha>0),\label{eq:gsingle}\\
 g(\omega,n)&=\omega+(n-1) &&(n\ge1),\label{eq:gomegan}\\
 g(\omega,\omega)&=\omega\cdot2,\label{eq:gomega}\\
 g(\omega+1,n)&=\omega+n &&(n\ge1),\label{eq:gomega1n}\\
 g(\omega+1,\omega)&=\omega\cdot2+1.\label{eq:gomega1omega}
\end{align}
Here is the verification of the steps involving limits. For
\eqref{eq:gomegan}, reducing the infinite coordinate gives finite values
cofinal in $\omega$, and reducing the finite coordinate gives the successors
$\omega+1,\ldots,\omega+(n-1)$. For \eqref{eq:gomega}, the values from
\eqref{eq:gomegan}, followed by $+1$, have supremum $\omega\cdot2$.
For \eqref{eq:gomega1n}, reducing $\omega+1$ to $\omega$ gives the value
$\omega+n$, and all other branches are bounded by it. Finally, in
\eqref{eq:gomega1omega}, reducing the first coordinate to $\omega$ gives
$g(\omega,\omega)+1=\omega\cdot2+1$, while reducing either coordinate to a
finite ordinal contributes values with supremum only $\omega\cdot2$.

\begin{proposition}[Abstract graph invariance fails in the infinite case]\label{prop:infinitegraph}
The countable wpos
\[
 X=\omega\sqcup\omega,\qquad Y=(\omega+1)\sqcup\omega
\]
have isomorphic incomparability graphs but different friendly order types:
\[
 \incg(X)\cong\incg(Y)\cong K_{\aleph_0,\aleph_0},\qquad
 \fot(X)=\omega\cdot2\ne\omega\cdot2+1=\fot(Y).
\]
\end{proposition}

\begin{proof}
Each side of either disjoint sum is a countably infinite chain; all and only
cross-chain pairs are incomparable. The graph assertion follows. The rank
values are \eqref{eq:gomega} and \eqref{eq:gomega1omega}.
\end{proof}

\begin{proposition}[Connectivity and maximal order type do not suffice]\label{prop:infinitemot}
Let $Z=(\omega\cdot2)\sqcup1$ and $Y=(\omega+1)\sqcup\omega$.
Both incomparability graphs are connected, and
\[
 \mot(Z)=\mot(Y)=\omega\cdot2+1,
\]
but
\[
 \fot(Z)=\omega\cdot2,\qquad \fot(Y)=\omega\cdot2+1.
\]
\end{proposition}

\begin{proof}
Connectivity is immediate from the cross-chain edges. The friendly values
follow from \eqref{eq:gsingle} and \eqref{eq:gomega1omega}.
For clarity, the maximal-order-type values can also be checked from bad-tree
ranks. Put $h(\alpha,\beta)=\mot(\alpha\sqcup\beta)$. It has the same
recursion as \eqref{eq:g}, but boundary values
$h(\alpha,0)=\alpha$ and $h(0,\beta)=\beta$. Induction gives
$h(\alpha,1)=\alpha+1$,
$h(\omega,n)=\omega+n$,
$h(\omega,\omega)=\omega\cdot2$, and
$h(\omega+1,n)=\omega+n+1$.
At $(\omega+1,\omega)$, the branch through $(\omega,\omega)$ contributes
$\omega\cdot2+1$; the remaining cofinal branches have supremum
$\omega\cdot2$. Thus $h(\omega+1,\omega)=\omega\cdot2+1$ as claimed.
\end{proof}

These examples put a precise limit on the proposed method. Decomposition
into ordered incomparability components still works for every wpo, and
Theorem~\ref{thm:transfinite} reduces the problem to individual components.
What fails is the replacement of an arbitrary infinite component by a
finite-style graph rank or a universal predecessor operation on its maximal
order type. Its internal well-founded order structure can matter.

\section{Algorithms and reproducible verification}

\subsection{Value computation and witness construction}

If the incomparability graph is supplied as an adjacency list with $n$
vertices and $m$ edges, graph traversal computes its number of components
and hence $\fot(P)$ in $O(n+m)$ graph operations. If a full comparison
matrix is supplied, building the incomparability graph takes $O(n^2)$
time. If only Hasse-cover relations are supplied, transitive closure must
be computed first; a straightforward implementation uses $O(n^3)$ time.
The graph-only value algorithm does not require a particular orientation
of the comparable pairs, but construction of a legal witness does.

The constructive proof supplies a polynomial-time witness algorithm:
find the ordered components, build a connected-prefix linear extension in
each, and read all but the first point of each extension in reverse order,
starting from the highest component. A direct comparison-matrix
implementation takes at most $O(n^3)$ time. More careful maintenance of
minimal candidates improves this bound, but is not needed for the claim.

The reference exact-rank computation is deliberately separate:
\begin{lstlisting}[language=Python]
@cache
def rank(mask):
    answer = 0
    for x in vertices(mask):
        if incomparable[x] & mask:
            next_mask = mask & ~upset[x]
            answer = max(answer, 1 + rank(next_mask))
    return answer
\end{lstlisting}
It uses neither connected components nor the proved formula. At most $2^n$
residual masks are examined, with at most $n$ candidate moves at each mask:
$O(n2^n)$ bit-mask operations and $O(2^n)$ stored ranks, before accounting
for the bit length of the masks. This is an exponential audit algorithm,
not the recommended value algorithm for large finite inputs.

The supplied implementation uses arbitrary-precision integer masks. The
complexity bounds just stated count graph/comparison operations or abstract
bit-mask operations; they should not be confused with constant-time bounds
on arbitrarily large Python integers.

\subsection{Exhaustive finite checks}

A poset on integer labels $0,\ldots,n-1$ is \emph{naturally labelled} here
when $x<_P y$ implies $x<y$ as integers. The generator adds the largest label
as a new maximal point, choosing any order ideal of the previous poset as
its predecessor set. This constructs every naturally labelled relation
exactly once: deleting the largest label uniquely recovers the earlier
poset and the chosen ideal.

Every finite poset admits a linear extension, so every isomorphism type has
at least one naturally labelled representative. The enumeration is not an
isomorph-free enumeration and is not a count of all arbitrary labellings.
This distinction is relevant to interpreting the following actual run.

\begin{center}
\begin{tabular}{@{}rrrr@{}}
\toprule
$n$ & Naturally labelled relations & Exact-rank comparisons & Result \\
\midrule
0 & 1 & 1 & pass\\
1 & 1 & 1 & pass\\
2 & 2 & 2 & pass\\
3 & 7 & 7 & pass\\
4 & 40 & 40 & pass\\
5 & 357 & 357 & pass\\
6 & 4,824 & 4,824 & pass\\
7 & 96,428 & 96,428 & pass\\
\midrule
Total & 101,660 & 101,660 & pass\\
\bottomrule
\end{tabular}
\end{center}

For every listed relation the audit compared the exact residual rank with
$n-c$, computed the incidence-matrix rank over $\F_2$, constructed an
optimal witness, and checked each selected point and friend against the
actual residual. A further 351 seeded relabelling tests checked that the
implementation does not accidentally depend on natural labels.

For additional reproducibility, the distributions by friendly order type
are stored in \path{data/rank_distributions.csv}. For example, at $n=7$
the numbers with friendly ranks $0,1,\ldots,6$ are
\[
 1,\ 6,\ 30,\ 160,\ 1047,\ 8906,\ 86278,
\]
respectively. These are outputs of the supplied enumeration, not asserted
new integer sequences or independent proof certificates for infinite cases.

\subsection{Product and substitution checks}

The product audit tested all ordered pairs of naturally labelled factors
of sizes two through five. There are $2+7+40+357=406$ such factor
relations, giving $406^2=164{,}836$ product cases. For each case it compared
the directly computed number of product incomparability components with
Theorem~\ref{thm:productgraph}. The 241 cases with product size at most nine
were also checked by the independent exact-rank recurrence. The other
product cases were \emph{graph-structure checks}, not exponential exact-rank
checks on up to 25 points.

The substitution audit used every naturally labelled index of sizes zero
through three and every assignment of nonempty naturally labelled fibers
of sizes one through three. There are ten available nonempty fiber
relations. The total is
\[
 1+1\cdot10+2\cdot10^2+7\cdot10^3=7{,}211.
\]
All 7,211 cases, each with at most nine points, passed both the exact-rank
comparison and the direct graph-rank comparison.

The recorded run used Python 3.13.5 and completed all audits in about 31
seconds in the execution environment used for this report. This timing is
not a portable benchmark or an asymptotic performance claim. Exact counts,
individual timings, and outputs are in \texttt{data/audit\_results.json}.

\subsection{How to reproduce the package}

The code needs only the Python standard library and is written for Python
3.10 or later. From the package root run
\begin{lstlisting}[language=bash]
python3 code/audit.py --max-n 7 --output data
\end{lstlisting}
The check through eight elements is optional and substantially slower:
\begin{lstlisting}[language=bash]
python3 code/audit.py --max-n 8 --output data_n8
\end{lstlisting}
No eight-element exhaustive run is claimed in this article. The finite
proof covers all sizes independently of these computations.

A simple library use is
\begin{lstlisting}[language=Python]
from friendly_order import chain, cartesian

P = cartesian(chain(2), chain(3))
assert P.friendly_formula() == 3
assert P.exact_friendly_rank() == 3
certificate = P.optimal_witness()
P.verify_witness(certificate)
\end{lstlisting}
Run this with \texttt{code/} on the Python module search path, for example
from inside that directory. The data file \texttt{examples.json} contains
explicit upset masks, component extensions, omitted roots, and friend
certificates for the examples in this article.

\Needspace{9\baselineskip}
\section{Research status, dependencies, and remaining questions}

\subsection{What has and has not been established}

The finite component theorem is a universal statement with a complete
constructive proof. The substitution and Cartesian formulas are consequences
of proved graph-structure lemmas. The well-ordered-sum formula is proved by
transfinite induction, not extrapolated from finite tests. The countable
counterexamples are verified by explicit ordinal recurrences.

The multiset-width conclusions have a different dependency: they apply
Vialard's published Theorem~\ref{thm:vialard}. The present note does not
replace its general proof. Nor does it solve the broader problem of
computing friendly order type on arbitrary infinite incomparability-connected
wpos. The finite-component class is a significant exact domain, but many
infinite wpos have infinite components and lie outside it.

The relation to other ordinal-invariant work is consistent with the
perspective of the finitary-powerset paper of Abriola and collaborators,
which emphasizes that additional invariants or restricted structural classes
may be needed to regain compositional computation \cite{abriola2024}.
Here, for finite bases, the additional information is an elementary graph
component count. In infinite examples even the abstract graph can lose the
necessary order information.

\subsection{Priority and literature check}

The source audit examined Vialard's published MFCS 2023 paper, the author's
2024 thesis copy, and the finitary-powerset follow-up. Searches included
``friendly order type'', ``friendly order type finite posets'', and
combinations with ``components'' and ``incomparability''. The published
conclusion and the thesis conclusion explicitly present further study of
friendly order type as a research direction. No finite component formula
was located in the materials inspected.

This establishes the provenance of the question and describes the search,
not a proof of novelty. The formula may occur elsewhere, be known
informally, or be obtainable from a result not found in that search. The
appropriate status is \emph{a proved result in this unrefereed manuscript,
with priority unverified}. The source lookup date was 19 September 2026.

\subsection{Precise directions left by the result}

The reduction to incomparability components suggests a focused remaining
problem: classify $\fot(C)$ for infinite incomparability-connected wpos $C$.
Propositions~\ref{prop:infinitegraph} and~\ref{prop:infinitemot} show that both
the abstract graph and maximal order type alone can be insufficient. A useful
next invariant would have to retain some interaction between finite
terminal structure and infinite bad-sequence ranks.

Another concrete question concerns \emph{which} subsets of a finite poset
can occur as the selected support of an open-ended bad sequence. The rank
theorem determines the maximum support size and constructs an optimum, but
does not characterize all feasible supports or identify their family with
the independent sets of a matroid. Confusing equality of maximal ranks with
an equality of combinatorial structures would be an unjustified extension.

Finally, efficient incremental witness maintenance can be studied alongside
Proposition~\ref{prop:insertion}. The numerical rank change under a valid
one-point insertion is completely determined by the old components met by
the new point. Updating a particular connected linear extension while
retaining a short certificate is a more detailed algorithmic question.

\section{Conclusion}

For finite posets, the residual definition of friendly order type conceals a
simple graph invariant:
\[
 \fot(P)=r(\incg(P))=|P|-\cc(P).
\]
The upper bound is a forest-rank potential, and the lower bound is a
connected-prefix linear-extension construction. Their combination gives an
exact value, an optimal sequence, and a spanning-forest certificate.

This answers the finite structural subproblem of the published program and
supplies further compositional formulas, notably
\[
 \fot(A\times B)=|A||B|-1-b(A)b(B)-t(A)t(B)
 \quad(|A|,|B|\ge2).
\]
Well-ordered additivity extends the evaluation to all wpos with finite
incomparability components, while explicit infinite examples explain why
additional structure is needed beyond that class. The reproducible checks
support the implementation and finite statements; the proofs, not the
checks, establish the general conclusions.

\appendix
\section{A compact proof certificate for the finite theorem}

The complete finite theorem can be checked against the following four
claims, each proved in the body of the article.
\begin{enumerate}[label=\textbf{\arabic*.},leftmargin=2em,itemsep=0.4em]
\item A legal move deletes $x$ but not a friend $y$. A spanning forest of
the residual can be extended by the edge $xy$, so graph rank drops by at
least one.
\item Incomparability components are linearly ordered ordinal-sum blocks.
The direction of comparison with an outside vertex cannot change along an
incomparability edge.
\item A connected finite component has a linear extension with connected
prefixes. Otherwise some nonempty proper ideal would be below every
remaining point and disconnect the graph.
\item Reverse each such extension, omit its first point, and process blocks
from high to low. Every selected point has an earlier-prefix friend. The
result has exactly one fewer selected point than vertices in each component.
\end{enumerate}
The first claim proves the bound $\fot(P)\le|P|-\cc(P)$, and the remaining
claims produce a witness attaining it. None assumes that every legal move
preserves connectivity, or that the omitted points must remain in the final
residual.

\section{Artifact manifest}

\begin{center}
\renewcommand{\arraystretch}{1.2}
\begin{tabular}{@{}p{0.41\textwidth}p{0.51\textwidth}@{}}
\toprule
File & Purpose \\
\midrule
\texttt{article.tex} & Complete source of this article; bibliography embedded.\\
\texttt{article.pdf} & Compiled article.\\
\texttt{README.md} & Build, reproduction, scope, and status instructions.\\
\texttt{build.sh} & Two-pass PDF build using \texttt{pdflatex}.\\
\texttt{code/friendly\_order.py} & Validated finite-poset representation, exact rank, graph formula, operations, and witnesses.\\
\texttt{code/audit.py} & Exhaustive and compositional verification driver.\\
\texttt{data/audit\_results.json} & Actual audit counts, timings, and pass status.\\
\texttt{data/audit\_console.txt} & Captured console output of the reported run.\\
\path{data/rank_distributions.csv} & Counts by size and friendly rank.\\
\texttt{data/examples.json} & Explicit finite examples and certificates.\\
\texttt{RESEARCH\_STATUS.md} & Claim/dependency boundary and source audit.\\
\bottomrule
\end{tabular}
\end{center}

\begin{thebibliography}{9}

\bibitem{vialard2023}
Isa Vialard.
\newblock Ordinal Measures of the Set of Finite Multisets.
\newblock In \emph{48th International Symposium on Mathematical Foundations
of Computer Science (MFCS 2023)}, LIPIcs \textbf{272}, article 87,
87:1--87:15, 2023.
\newblock DOI: \href{https://doi.org/10.4230/LIPIcs.MFCS.2023.87}
{10.4230/LIPIcs.MFCS.2023.87}.
\newblock \url{https://drops.dagstuhl.de/entities/document/10.4230/LIPIcs.MFCS.2023.87}.
\newblock In particular, Definition 3.3, Theorem 3.4, Proposition 4.1, and
the conclusion. Published version, rather than the early arXiv terminology.

\bibitem{vialard2024}
Isa Vialard.
\newblock \emph{Measuring Well Quasi-Orders and Complexity of Verification}.
\newblock Doctoral thesis, Universit\'e Paris-Saclay, 2024.
\newblock Author-hosted copy: \url{https://isavialard.github.io/home/mwqo.pdf}.
\newblock In particular, Sections 5.3--5.4 and the conclusion, printed
pp.~109--110. The page references here concern this 124-page author copy.

\bibitem{abriola2024}
Sergio Abriola, Simon Halfon, Aliaume Lopez, Sylvain Schmitz,
Philippe Schnoebelen, and Isa Vialard.
\newblock Measuring well-quasi-ordered finitary powersets.
\newblock arXiv:2312.14587, version 2, 17 July 2024.
\newblock \url{https://arxiv.org/abs/2312.14587}.
\newblock Used for research context, not as a dependency of the component
formula proved here.

\end{thebibliography}
\end{document}
